\begin{abstract}
Molecular property prediction benchmarks are commonly interpreted through aggregate model scores, yet those scores can be strongly shaped by how molecules are divided into training and test sets. Random splits may create chemically similar train/test pairs and therefore optimistic estimates, whereas scaffold splits are intended to test generalization to new molecular frameworks. However, scaffold splitting can also introduce target distribution shift and scaffold concentration, making it unclear whether a performance drop reflects chemical generalization difficulty or artifacts of the split itself. This study presents a reproducible split-diagnostic audit of six public molecular property datasets spanning classification and regression tasks. Across random, ordinary scaffold, and target-balanced scaffold protocols, fingerprint-based classical baselines are evaluated over five seeds. The resulting benchmarks are assessed through generalization gaps, target shifts, scaffold statistics, test-to-train Tanimoto similarity, outlier cases, and model-ranking stability. Random splits generally yield higher test-to-train chemical similarity and more optimistic performance estimates. Ordinary scaffold splits remove train/test scaffold overlap, but can create large target shifts in scaffold-concentrated datasets. Target-balanced scaffold splitting reduces target shift in several cases, but residual gaps and ranking changes remain dataset- and model-dependent. These results support reporting split diagnostics alongside model performance, because split choice can alter both the apparent difficulty of a dataset and the apparent superiority of a model.
\end{abstract}

\textbf{Keywords:} AI-driven drug discovery; molecular property prediction; scaffold split; benchmark audit; Tanimoto similarity; target distribution shift; generalization gap.
